\documentclass[11pt]{article}
\usepackage[a4paper,margin=2.5cm]{geometry}
\usepackage{amsmath, amssymb}
\usepackage{graphicx}
\usepackage{hyperref}
\usepackage{titlesec}
\usepackage{enumitem}

\title{\textbf{Tensor Omics: Semantic Geometry in Multi-Omics Analysis}}
\author{Bobito}
\date{}

\begin{document}
\maketitle

\section{Overview}

\textit{Tensor Omics} is a conceptual and computational framework for representing, analyzing, and integrating high-dimensional omics data—transcriptomics, proteomics, metabolomics, metagenomics, and genome-wide phenotype associations—using tensor-based geometry. Instead of treating genes or features in isolation, Tensor Omics places their quantification vectors into a shared semantic space defined by experimental conditions or biological contexts (e.g., tissues, timepoints, treatments). This space is used to construct geometric fields of expression differences, divergence vectors, and local second-order tensor structures, enabling both intuitive interpretation and advanced computational comparison.

By embedding multi-omics data into a unified or interoperable coordinate grid and comparing vectorial behaviors across dimensions and biological conditions, Tensor Omics supports fine-grained questions about gene family evolution, tissue specificity, functional divergence, microbial ecosystem dynamics, and integrated omics response profiling. The core of the method lies in its commitment to semantic geometry: each axis of the space carries biological meaning, each vector a measurable state, and each field a dynamic, spatially structured narrative of molecular function.

Tensor Omics offers native compatibility with vector-based AI methods, fast geometric comparison tools, and a fully reproducible, versionable implementation stack. It scales to transcriptomics of a single species, integrates proteomics and metabolomics across timepoints, and extends to metagenomic ecosystems and phenotype-genome association landscapes. From expression-based neofunctionalization detection to high-resolution microbiome response profiling, Tensor Omics transforms geometry into biology.

\section{1. Methods}

\subsection{1.1 Input Data and Initial Representation}

We assume as input an omics dataset—typically transcriptomics—with the following components:

\begin{itemize}[leftmargin=1.5em]
    \item A matrix of expression values $X \in \mathbb{R}_{\geq 0}^{n \times m}$, where $n$ is the number of genes (or features) and $m$ the number of experimental conditions (e.g., tissues, timepoints, treatments).
    \item Associated metadata for each gene and each sample, stored in structured tables or hash maps for indexing, grouping, and filtering.
\end{itemize}

Each row vector $\mathbf{x}_i \in \mathbb{R}^m$ represents the expression profile of gene $i$ across $m$ semantic dimensions (samples), forming the basis for all further geometric interpretation.

\subsection{1.2 Normalization}

To enable meaningful comparisons across genes and omics layers, the following normalization pipeline is applied:

\begin{enumerate}[label=(\roman*), leftmargin=1.5em]
    \item \textbf{Standard Deviation Scaling:} Each vector $\mathbf{x}_i$ is scaled by the standard deviation of its components:
    \[
        \mathbf{x}_i' = \frac{\mathbf{x}_i}{\mathrm{sd}(\mathbf{x}_i)}
    \]
    \item \textbf{Quantile Normalization:} Aimed at ensuring consistent expression distributions across samples. This is optional and depends on the experimental design; it is more suitable within omics types than across them.
    \item \textbf{Log Transformation:} To reduce the skew caused by extreme values and emphasize relative changes, we apply a log transformation:
    \[
        \mathbf{x}_i'' = \log(1 + \mathbf{x}_i')
    \]
\end{enumerate}

After normalization, each gene is now represented as a vector in $\mathbb{R}^m$ with a geometry that is interpretable across all $m$ experimental conditions. These normalized vectors define the primary semantic expression space.

\subsection{1.3 Construction of the Gene Expression Space}

The normalized expression vectors $\{\mathbf{x}_i''\}_{i=1}^n$ are embedded into a semantic coordinate space $\mathbb{R}^m$, where each axis corresponds to a biological dimension such as a tissue type, timepoint, treatment condition, or spatial location. This space—called the \textit{omics space}—forms the geometric domain for all downstream analysis.

Let $G$ denote a discrete grid imposed on $\mathbb{R}^m$. The grid resolution is determined automatically based on data density using a robust, reproducible strategy—e.g., the $0.8$ quantile of nearest-neighbor distances between vectors. The grid provides a canonical lattice to anchor vector fields, allowing field interpolation, modular detection, and multi-omics comparison.

\subsection{1.4 Interpolation and Smoothing}

Because real data is sparse and noisy, Tensor Omics applies interpolation and smoothing at several stages to enable robust field construction and biological interpretation:

\begin{itemize}[leftmargin=1.5em]
    \item \textbf{Stage 1:} After embedding the expression vectors into the space to stabilize geometry.
    \item \textbf{Stage 2:} After computing first-order tensor fields (e.g., flow directions) to clarify vector transitions.
    \item \textbf{Stage 3:} After identifying local modules (via Bobito or similar tools), to refine trace lines and surface curvature estimates.
\end{itemize}

Smoothing is implemented via kernel-weighted interpolation on the grid. At a grid point $\mathbf{g}$, the interpolated vector $\mathbf{v}(\mathbf{g})$ is computed as:
\[
    \mathbf{v}(\mathbf{g}) = \frac{\sum_{i} K(\|\mathbf{x}_i - \mathbf{g}\|) \cdot \mathbf{x}_i}{\sum_{i} K(\|\mathbf{x}_i - \mathbf{g}\|)}
\]
where $K(\cdot)$ is a kernel function (e.g., Gaussian) with a radius tied to the grid resolution. This approach enables local continuity without imposing global structure, retaining flexibility and interpretability.

\section{Family-Level Geometry and Outlier Detection}

\subsection{2.1 Gene Family Centroids}

Genes are grouped into families based on shared sequence, function, or evolutionary history. For a given gene family $F = \{ \mathbf{x}_1, \ldots, \mathbf{x}_k \}$, we define the centroid vector $\mathbf{c}_F$ as:

\[
    \mathbf{c}_F = \frac{1}{k} \sum_{i=1}^{k} \mathbf{x}_i
\]

This centroid represents the average expression behavior of the family and serves as a neutral reference point. It can optionally be replaced by the centroid of orthologs in a reference species or population when available.

\subsection{2.2 Difference Vectors}

For each gene $\mathbf{x}_i \in F$, we define its \textit{difference vector} $\mathbf{d}_i$ as the vector from the family centroid to the gene:

\[
    \mathbf{d}_i = \mathbf{x}_i - \mathbf{c}_F
\]

This difference vector captures the gene’s divergence in expression space relative to its family norm. It represents both direction (e.g., toward specific tissues) and magnitude (degree of divergence).

\subsection{2.3 Relative Distance Index (RDI)}

To distinguish meaningful divergence from noise, Tensor Omics uses a native empirical significance measure called the \textit{Relative Distance Index (RDI)}.

\textbf{Step 1:} Construct a background distribution of difference vector lengths $\|\mathbf{d}_j\|$ for all genes in all families.

\textbf{Step 2:} For a given gene $i$, its RDI is defined as:

\[
    \text{RDI}_i = \frac{\text{rank}(\|\mathbf{d}_i\|)}{\text{total number of distances}}
\]

This yields a value between 0 and 1 representing the gene’s relative divergence.

\textbf{Thresholding:} Genes with $\text{RDI} > 0.95$ are considered significant outliers—i.e., likely to have undergone expression divergence due to subfunctionalization, neofunctionalization, or other adaptive processes.

The RDI measure is:
\begin{itemize}[leftmargin=1.5em]
    \item \textit{Non-parametric} and data-driven.
    \item \textit{Applicable} across experiments without requiring a model.
    \item \textit{Composable} with metadata filters (e.g., tissues, conditions, species).
\end{itemize}

\section{Semantic Axes: Versatility, Specificity, and Divergence toward Tissue Identity}

\subsection{3.1 The Space Diagonal and Expression Neutrality}

Let the semantic axes of the expression space be denoted by $T = \{t_1, \ldots, t_n\}$ (e.g., tissues or conditions). In this $n$-dimensional space, the \textit{space diagonal} is the vector:

\[
    \mathbf{d}_{\text{diag}} = (1, 1, \ldots, 1)
\]

It represents an equal expression in all axes — a uniform, versatile expression profile. Any deviation from this vector indicates semantic bias (i.e., tissue or condition specificity).

\subsection{3.2 Angular Deviation from the Diagonal}

For a gene expression vector $\mathbf{x}_i$, we compute the angle $\theta_i$ to the space diagonal using the cosine similarity:

\[
    \theta_i = \arccos\left( \frac{ \mathbf{x}_i \cdot \mathbf{d}_{\text{diag}} }{ \|\mathbf{x}_i\| \cdot \|\mathbf{d}_{\text{diag}}\| } \right)
\]

\textbf{Interpretation:}
\begin{itemize}[leftmargin=1.5em]
    \item $\theta \approx 0^\circ$: Uniform expression (tissue-versatile gene).
    \item $\theta \gg 0^\circ$: Biased expression (tissue-specific gene).
\end{itemize}

This measure allows rapid geometric identification of genes that have evolved specialization, for instance after duplication. Tracking these angles across orthologs and paralogs enables inference of sub- and neofunctionalization dynamics.

\subsection{3.3 Practical Applications}

\textbf{Use cases:}
\begin{itemize}[leftmargin=1.5em]
    \item \textit{Trichoplax stress response (Heigwer project)}: Are stress-inducible genes more tissue-specific under stress? Does angular deviation increase?
    \item \textit{Chicken microbiome (Dusel project)}: Do certain taxa exhibit more specialized expression under different diets?
    \item \textit{Wound healing (Troidl project)}: Which genes lose or gain versatility across genotypes or treatments?
\end{itemize}

The angle $\theta$ is compact, native, and suitable for large-scale screening.

\section{4. Relative Axes Expression: Clock Plots and Semantic Rotation}

\subsection{4.1 Projection onto the Diagonal-Orthogonal Plane}

To study shifts in expression directionality—e.g., between paralogs or across time—we define a projection onto the plane orthogonal to the space diagonal $\mathbf{d}_{\text{diag}}$.

Let $\mathbf{x} \in \mathbb{R}^n$ be any expression vector. Its projection $\mathbf{x}^\perp$ onto the diagonal-orthogonal plane is given by:

\[
\mathbf{x}^\perp = \mathbf{x} - \left( \frac{\mathbf{x} \cdot \mathbf{d}_{\text{diag}}}{\|\mathbf{d}_{\text{diag}}\|^2} \right) \mathbf{d}_{\text{diag}}
\]

This yields the component of $\mathbf{x}$ that lies in the plane where tissue biases can rotate without affecting total expression magnitude.

\subsection{4.2 Clock Hand Vectors and Angle $\phi$}

Let:
\begin{itemize}[leftmargin=1.5em]
    \item $\mathbf{o}^\perp$ be the projection of the ortholog expression vector onto the diagonal-orthogonal plane.
    \item $\mathbf{p}^\perp$ be the projection of the paralog expression vector.
    \item The \textbf{clock hand vectors} are:
    \[
    \mathbf{v}_o = \mathbf{o}^\perp - \mathbf{D}, \quad \mathbf{v}_p = \mathbf{p}^\perp - \mathbf{D}
    \]
    where $\mathbf{D}$ is the diagonal point in the plane used as the clock center (e.g., mean projection).
\end{itemize}

The angle $\phi$ between clock hands:

\[
\phi = \arccos \left( \frac{ \mathbf{v}_o \cdot \mathbf{v}_p }{ \|\mathbf{v}_o\| \cdot \|\mathbf{v}_p\| } \right)
\]

\textbf{Interpretation:}
\begin{itemize}[leftmargin=1.5em]
    \item $\phi$ reflects directional change in tissue bias.
    \item $\phi = 0$: perfect conservation of semantic orientation.
    \item $\phi > 0$: divergence in tissue preference.
\end{itemize}

\subsection{4.3 Difference Vector and Unit Difference Vector}

Define the \textbf{difference vector}:
\[
\mathbf{d}_{\text{raw}} = \mathbf{p}^\perp - \mathbf{o}^\perp
\]

And the \textbf{unit difference vector}:
\[
\hat{\mathbf{d}} = \frac{ \mathbf{d}_{\text{raw}} }{ \| \mathbf{d}_{\text{raw}} \| }
\]

These vectors encode the direction and magnitude of tissue expression shift between two expression states.

\subsection{4.4 Biological Interpretations and Use Cases}

\begin{itemize}[leftmargin=1.5em]
    \item \textbf{Neofunctionalization:} Detected as a large $\phi$ and strong $\mathbf{d}_{\text{raw}}$ pointing into a new tissue direction.
    \item \textbf{Subfunctionalization:} Identified when paralogs occupy opposite sectors of the clock and their vector sum approximates the ortholog.
    \item \textbf{Time Series Clock Plots:} If multiple timepoints exist (e.g., in wound healing), the angle $\phi(t)$ can be plotted to show semantic rotation over time.
\end{itemize}

This allows visualization and interpretation of evolving expression trajectories as rotations in semantic space.

\section{5. Meta-Data Filtering and Comparative Analysis}

\subsection{5.1 Subsetting the Semantic Vector Space}

The Tensor Omics framework supports meta-data–driven filtering of expression vectors. This enables targeted analysis of semantic flow and divergence within biologically meaningful subsets of the data.

Let:
\begin{itemize}[leftmargin=1.5em]
    \item $\mathcal{M}$ be a meta-data table with $n$ rows (e.g., one per gene) and $m$ columns (e.g., taxon, gene family, function, etc.).
    \item Subsets of expression vectors can be selected by applying filters over $\mathcal{M}$, e.g.:
    \[
    \text{select: taxon} = \text{``Arabidopsis''} \land \text{GO term} = \text{``photosynthesis''}
    \]
\end{itemize}

\subsection{5.2 Distributional Comparison of Expression Differences}

Once a semantic subset is selected, difference vectors (e.g., between paralog and ortholog) can be compared across groups using statistical methods:

\begin{itemize}[leftmargin=1.5em]
    \item \textbf{Magnitude Comparisons:} Compare the Euclidean lengths of difference vectors.
    \item \textbf{Angular Comparisons:} Compare the directionality of normalized difference vectors using circular statistics (e.g., Watson–Williams test or angular variance).
    \item \textbf{Vector Field Comparisons:} Compare first-order tensor fields (e.g., net mean flow) across species or conditions.
\end{itemize}

Example:
\begin{quote}
Compare the tissue shift direction (normalized $\hat{\mathbf{d}}$) for genes annotated as photosynthesis-related between \textit{Arabidopsis} and \textit{Camelina}. Use the Wilcoxon rank-sum test to determine if shift magnitudes differ significantly.
\end{quote}

\subsection{5.3 Interpretation of Semantic Flow}

The aggregate vector field (mean or weighted mean of difference vectors in a subset) yields interpretable biological insight:
\begin{itemize}[leftmargin=1.5em]
    \item \textbf{Net Mean Direction:} Dominant tissue reprogramming trend in a functional group.
    \item \textbf{Flow Strength:} Average magnitude of functional divergence.
    \item \textbf{Direction Clustering:} Tight angular distribution suggests concerted biological response; diffuse pattern implies functional drift.
\end{itemize}

\subsection{5.4 Example Applications from the Proposal}

\begin{itemize}[leftmargin=1.5em]
    \item \textbf{Microbiome under dietary shift:} Compare gene expression divergence vectors for microbial taxa under different diets. Highlight functional adaptation or specialization.
    \item \textbf{Single-cell stress response:} Filter by transcription factors and assess whether neofunctionalization occurs under heat or oxidative stress.
    \item \textbf{Wound healing:} Use time-trajectory vector fields to find gene modules with early versus late semantic activation.
\end{itemize}

\section{6. Software Architecture and Implementation Plan}

\subsection{6.1 Language and Framework Choice}

The Tensor Omics implementation is deliberately minimal, efficient, and reproducible by design. The core computation stack consists of:

\begin{itemize}[leftmargin=1.5em]
    \item \textbf{R:} Primary environment for data ingestion (e.g., via \texttt{tximport}), scripting, statistical analysis, and interaction with the Fortran backend.
    \item \textbf{Fortran:} High-performance engine for tensor field calculations, vector operations, serialization, and WebAssembly (WASM) compilation.
    \item \textbf{LAPACK:} Used via \texttt{DGESDD} for extracting dominant eigenvectors of second-order tensors (surface curvature).
    \item \textbf{FortranContainers:} Provides hash maps and dynamic lists, used for metadata mapping and grid storage.
\end{itemize}

Python is explicitly avoided. If needed, specific functions may be implemented using Rcpp with C++ libraries, maintaining a modular and lightweight R package.

\subsection{6.2 Data Structure and Serialization}

All omics data, vectors, and tensor fields are stored in a unified binary format called \texttt{.txdata}, containing:

\begin{itemize}[leftmargin=1.5em]
    \item \texttt{axes}: Character array of semantic dimensions (e.g., tissue types).
    \item \texttt{omics\_vecs}: Matrix of real numbers, each column a high-dimensional expression vector.
    \item \texttt{meta}: Consists of:
    \begin{itemize}
        \item Flat arrays (\texttt{real\_data}, \texttt{int\_data}, \texttt{char\_data}) to store metadata values.
        \item \texttt{MetaColumn} structures to index and describe each metadata column.
        \item Hash maps to map identifiers (e.g., gene ID to index).
    \end{itemize}
    \item \texttt{tensor\_field\_frst\_ord}: First-order tensor field (e.g., directional divergence vectors).
    \item \texttt{tensor\_field\_scnd\_ord}: Second-order tensor field (e.g., local curvature matrices).
\end{itemize}

Serialization includes a self-descriptive header, followed by raw binary blocks. All structures are designed for compatibility with WASM (no dynamic allocation, no co-arrays, no nested structures).

\subsection{6.3 Workbench for Interactive Visualization}

The Tensor Omics Workbench is a browser-based GUI, built with:

\begin{itemize}[leftmargin=1.5em]
    \item \textbf{HTML, CSS, JavaScript}
    \item \textbf{D3.js} for interactive 2D/3D plots.
    \item \textbf{WebAssembly} (compiled from Fortran) for fast projection and tensor lookup.
\end{itemize}

Key features include:
\begin{itemize}[leftmargin=1.5em]
    \item Zooming, axis toggling, and spatial rotation.
    \item Meta-data search and vector coloring (by category, value, or symbol).
    \item Projection into planes, clock plot rendering, and grid-level inspection.
    \item Export of plot state (coordinates, zoom, coloring, axis configuration) for reproducibility and sharing.
\end{itemize}

\subsection{6.4 Patching and Reproducibility}

All changes to the data can be serialized as lightweight binary patches. These patches:
\begin{itemize}[leftmargin=1.5em]
    \item Are logged with a timestamp and brief description.
    \item Can be streamed from the analysis server (e.g., via VSCode over SSH) to the browser.
    \item Are versionable using Git or DVC and can be included in reproducible lab notebooks (e.g., RMarkdown).
\end{itemize}

This infrastructure supports iterative and fully traceable multi-omics analysis pipelines.

\section{7. Summary and Evaluation of Tensor Omics}

Tensor Omics is a novel framework for the geometric and semantic analysis of high-dimensional omics data. At its core lies a deceptively simple idea: omics data, especially from multi-condition or multi-tissue experiments, naturally lives in vector spaces. Instead of abstracting or reducing this complexity, Tensor Omics embraces it—turning raw data into mathematically structured vector and tensor fields. Each axis retains biological meaning, and operations like interpolation, smoothing, and derivative computation are implemented natively in this semantically grounded space.

This approach allows for new forms of biological insight:

\begin{itemize}[leftmargin=1.5em]
    \item \textbf{Vector Divergence and Flow:} Divergence vectors reveal functional shifts across gene families.
    \item \textbf{Clock Plots and Angular Analysis:} Trajectories in relative axis expression space reveal neofunctionalization and tissue specificity.
    \item \textbf{Tensor Fields and Curvature:} Local curvature and stress reveal module surfaces and response geometry.
    \item \textbf{Integration Across Omics:} With bijective semantic mapping and unified grid design, transcriptomics, proteomics, and metabolomics can be compared geometrically.
\end{itemize}

Unlike most current bioinformatics tools, Tensor Omics makes minimal statistical assumptions. Instead, it provides a platform—semantic geometry—for scientists to formulate, test, and visualize hypotheses natively. The system's architecture (Fortran + R + WASM + D3) was chosen for its robustness, reproducibility, and clarity.

\textbf{Why this matters:} Many breakthroughs in genomics and systems biology today are limited by the absence of tools that represent the full semantic and spatial structure of omics data. Tensor Omics proposes a path forward, not by adding layers of abstraction, but by grounding analysis in the data's intrinsic geometry. It transforms tables of numbers into living mathematical fields—into maps of biological processes in motion.

\textbf{Tensor Omics is the toolset. Semantic Geometry is the worldview.}

\end{document}
